%قبل از شروع مسئله ترجیحاً چیزی ننویسید. اما اگر خواستید بنویسید، در قالب لیستی (فقرات و شمارش) بنویسید.
%برای اضافه کردن نمره، بعدی از شروع‌مسئله، [x] را اضافه کنید و به جای x نمره مورد نظر را قرار دهید. برای زیرمسئله‌ها هم به همین صورت است.
%\شروع‌مسئله[10]
\شروع‌مسئله 
در این مرحله شما باید الگوریتمی برای تشخیص برخورد دو چندضلعی پیاده‌سازی کنید. برای این مسأله قسمت \lr{pairwise pruning} 
\href{https://en.wikipedia.org/wiki/Collision_detection}{این لینک}
را مطالعه کنید.
 دو چندضلعی به صورت دو لیست از مثلث‌ها به شما داده می‌شود.(چندضلعی‌ها مثلث‌بندی شده‌اند) این دو چندضلعی در صورتی برخورد دارند که یکی از مثلث‌هایشان با هم برخورد داشته باشند. پس شما باید الگوریتمی را پیاده‌سازی کنید که برخورد بین مثلث‌های این دو چندضلعی را تشخیص دهد. برای پیاده‌سازی بهینۀ این مسأله باید از الگوریتم‌های \lr{bounding volume hierarchy} استفاده کنید. در مورد انواع این الگوریتم‌ها جست‌و‌جو کنید. 
\href{http://www-ljk.imag.fr/Publications/Basilic/com.lmc.publi.PUBLI_Inproceedings@117681e94b6_1860ffd/bounding_volume_hierarchies.pdf}{این لینک}
و 
\href{https://cg.informatik.uni-freiburg.de/course_notes/sim_06_boundingVolumeHierarchies.pdf}{این لینک}
 می‌توانند به شناخت انواع این الگوریتم‌ها کمک کنند. پیشنهاد می‌شود برای سادگی و درگیر نشدن با قسمت هندسی این الگوریتم از \lr{AABB} در پیاده‌سازی این الگوریتم استفاده کنید. هر چند در صورتی که الگوریتم‌های پیچیده‌تری را پیاده‌سازی کنید تست‌کیس‌های بیشتری را پاسخ خواهید داد. و در صورتی هم که به صورت \lr{brute force} عمل کنید به دلیل سختی تست‌کیس‌ها نمرۀ پایینی خواهید گرفت. در بررسی برخورد دو مثلث با یکدیگر می‌توانید از کدهای موجود در اینترنت استفاده کنید. این تنها جایی است که چنین امکانی دارید و در صورت مشاهدۀ این مورد در هر جای دیگر از کد نمره‌ای به شما تعلق نخواهد گرفت. تنها از زبان برنامه‌نویسی پایتون می‌توانید استفاده کنید و امکان استفاده از هیج کتابخانه‌ای را ندارید. همچنین امکان استفاده از روش‌های \lr{multithreading} وجود ندارد و برنامه‌تان باید به صورت ترتیبی اجرا شود. به ازای یک ورودی برنامه‌تان همیشه باید جواب ثابتی تولید کند. کد شما باید کاملاً درست کار کند. به این معنا که در صورتی که کدتان \lr{time limit} نخورد باید پاسخ درستی بدهد. پیشنهاد می‌شود که جست‌و‌جو در مورد الگوریتم‌های مختلف این مسأله را به روزهای آخر موکول نکنید چرا که فرصت پیاده‌سازی، تست و بهبود آن‌ها را پیدا نخواهید کرد.
 
 \متن‌سیاه{ورودی و خروجی:}
 در هر تست‌کیس ۳۰ جفت چندضلعی بررسی می‌شود. ورودی زیر ۳۰ بار به برنامۀ شما داده می‌شود:
 \\
 ابتدا مثلث‌های چندضلعی اول به عنوان ورودی داده می‌شوند. در هر خط یک مثلث به صورت 
 $x_1, y_1, x_2, y_2, x_3, y_3$
  داده می‌شود. پس از پایان مثلث‌های چندضلعی اول \lr{end1} می‌آید و سپس مثلث‌های چندضلعی دوم به عنوان ورودی داده می‌شوند. در آخر نیز دستور \lr{end2} می‌آید. در صورتی که دو چندضلعی برخورد داشتند عدد ۱ و در غیر این صورت عدد ۰ در خروجی چاپ می‌شود. در نتیجه خروجی شما باید شامل ۳۰ خط از اعداد ۰ و ۱ باشد. در ادامه دو نمونه آورده شده است که یک جفت چندضلعی را بررسی می‌کنند.
 
 \صفحه‌جدید
  نمونۀ اول:
  
  %\lstinputlisting{source_filename.py}
  \begin{latin}
  \begin{lstlisting}
  # input
  1 1 1 2 2 1
  1 2 2 1 3 3
  end1
  1 3 2 3 3 1
  end2
  # output
  1
  \end{lstlisting}
  \end{latin}
 نمونۀ دوم:

%\lstinputlisting{source_filename.py}
  \begin{latin}
  \begin{lstlisting}
  # input
  1 1 1 3 2 3
  end1
  1 1 2 3 3 1
  end2
  # output
  0
  \end{lstlisting}
  \end{latin}

%\متن‌سیاه{سوال عملی اول:}\\
%\متن‌سیاه{دراپ‌تست}\\

%\متن‌سیاه{ورودی:}\\
%در خط اول ورودی، ابتدا k و سپس m داده می‌شود.\\
%\متن‌سیاه{خروجی:}\\
%در تنها سطر خروجی، تعداد آیفون‌های مورد نیاز می‌آید. اگر در هیچ صورتی چنین کاری ممکن نبود ۱- چاپ می‌شود.\\
%\متن‌سیاه{محدودیت‌ها:}\\
%$1\leq k\leq 10^4$

%$m\leq10^9$
%\\
%\textbf{\متن‌سیاه{ورودی و خروجی نمونه}\\
%% تست اول
%\begin{latin}
% \begin{tabularx}{\textwidth}{|X | X|}
% \hline
% \textbf{stdin} &
% \textbf{stdout}\\
% \hline
% 100 100&
% 1\\
% \hline
% 
% \end{tabularx}
%\end{latin}
%
%% تست دوم
%
%\begin{latin}
% \begin{tabularx}{\textwidth}{|X | X|}
% \hline
% \textbf{stdin} &
% \textbf{stdout}\\
% \hline
% 340 20&
% 3\\
% \hline
% 
% \end{tabularx}
%\end{latin}}
% صورت اصلی مسئله را در این‌جا بنویسید.
%اگر مسئله چند قسمت داشت به این صورت عمل کنید:
%\شروع{زیرمسائل}
%\زیرمسئله[0] % مسئلۀ ۱آ
%\زیرمسئله % مسئلۀ ۱ب
%\زیرمسئله % مسئلۀ ۱پ
%\زیرمسئله % مسئلۀ ۱ت
%\پایان{زیرمسائل}




%%%%%%%%%%%%%%%%%%%%%%%%%%%%%%%%%%%%%%%%%%%%%%%%%%%%%
% برخی دستورات مفید برای تمرین‌ها. %
%%%%%%%%%%%%%%%%%%%%%%%%%%%%%%%%%%%%%%%%%%%%%%%%%%%%%

% فرمول‌های نمونه:
%$f(x) = \Theta(x)$
%$T(x, y) \leq \log(x) + 2^y + \binom{2n}{n}$
% {\tt code\_function}


% لیست موارد بدون شماره‌گذاری:
%\شروع{فقرات}
% \فقره مورد اول 
% \شروع{فقرات}
% \فقره مورد اول داخلی 
% \شروع{فقرات}
% \فقره اولین مورد خیلی داخلی 
% \فقره دومین مورد خیلی داخلی 
% \پایان{فقرات}
% \فقره مورد دوم داخلی 
% \پایان{فقرات}
% \فقره مورد دوم 
%\پایان{فقرات}

% لیست موارد شماره‌گذاری شده:
%\شروع{شمارش}
% \فقره یک مورد 
% \فقره یک مورد دیگر 
% \فقره باز هم یک مورد دیگر
%\پایان{شمارش}

% برای نوشتن یک فرمول چند خطی که یک شماره بخوره. برای حالت‌های دیگر به سایت زیر مراجعه کنید:
% https://www.overleaf.com/learn/latex/Aligning_equations_with_amsmath
%\begin{align} 
% \begin{split}
% (x+y)^3 &= (x+y)^2(x+y) \\
% &= (x^2+2xy+y^2)(x+y) \\
% &= (x^3+2x^2y+xy^2) + (x^2y+2xy^2+y^3) \\
% &= x^3+3x^2y+3xy^2+y^3
% \end{split} 
%\end{align}

% نوشتن ماتریس و آرایه دوبعدی
%\begin{align}
% A = 
% \begin{bmatrix}
% A_{11} & A_{21} \\
% A_{21} & A_{22}
% \end{bmatrix}
%\end{align}

% اضافه کردن تصویر با شرح. برای اضافه کردن تصویر بدون شرح، فقط دستور درج‌تصویر را استفاده کنید.
%\شروع{شکل}[h!]
%\تنظیم‌ازوسط
% \درج‌تصویر[پهنا=0.5\پهنای‌متن]{img.jpg}
%\شرح{شرح تصویر}
%\برچسب{تصویر:۱}
%\پایان{شکل}
%\پایان{زیرمسائل}

%کد زیر را ببینید.
%%\lstinputlisting{source_filename.py}
%\begin{latin}
%\begin{lstlisting}
%# Import the modules
%import sys
%import random
%
%ans = True
%
%while ans:
%question = raw_input("Ask the magic 8 ball a question: (press enter to quit) ")
%
%answers = random.randint(1,8)
%
%if question == "":
%sys.exit()
%
%elif answers == 1:
%print "It is certain"
%
%elif answers == 2:
%print "Outlook good"
%
%elif answers == 3:
%print "You may rely on it"
%
%elif answers == 4:
%print "Ask again later"
%
%elif answers == 5:
%print "Concentrate and ask again"
%
%elif answers == 6:
%print "Reply hazy, try again"
%
%elif answers == 7:
%print "My reply is no"
%
%elif answers == 8:
%print "My sources say no"
%\end{lstlisting}
%\end{latin}